\documentclass[11pt,a4paper]{article}
\usepackage{cmap}
\usepackage[T2A]{fontenc}
\usepackage[utf8]{inputenc}
\usepackage[russian]{babel}
\usepackage{DejaVuSerif}
\usepackage{amsmath,amssymb,amsthm}
\usepackage{geometry}
\geometry{left=20mm,right=20mm,top=20mm,bottom=22mm}
\usepackage{tikz}
\usetikzlibrary{arrows.meta,positioning,shapes.geometric}
\usepackage{xcolor}
\usepackage{tabularx,array,booktabs}
\usepackage{fancyhdr}
\usepackage{enumitem}
\usepackage{microtype}
\usepackage{needspace}
\usepackage{setspace}
\usepackage{hyperref}
\hypersetup{hidelinks,pdftitle={Простые и составные числа},pdfauthor={Учебно-методическое пособие}}

\definecolor{ink}{HTML}{263746}
\definecolor{navy}{HTML}{304E69}
\definecolor{slate}{HTML}{758A9B}
\definecolor{muted}{HTML}{5F6E79}
\definecolor{pale}{HTML}{E7EDF1}
\color{ink}
\setstretch{1.09}
\setlength{\parindent}{0pt}
\setlength{\parskip}{4.5pt}
\setlist[itemize]{leftmargin=6mm,itemsep=2pt,topsep=3pt}
\setlist[enumerate]{leftmargin=8mm,itemsep=2pt,topsep=3pt}
\setcounter{secnumdepth}{0}
\pagestyle{fancy}
\fancyhf{}
\fancyhead[L]{\footnotesize\color{muted} ЕГЭ по математике\,|\,Целые числа}
\fancyhead[R]{\footnotesize\color{muted} Простые и составные числа}
\fancyfoot[C]{\small\color{muted}\thepage}
\renewcommand{\headrulewidth}{0.35pt}
\setlength{\headheight}{14pt}
\renewcommand{\arraystretch}{1.37}

\newcommand{\blocktitle}[2]{\clearpage\section{Блок #1. #2}\vspace{-2mm}\noindent\color{slate}\rule{\linewidth}{.6pt}\color{ink}\vspace{2mm}}
\newcommand{\concept}[1]{\par\needspace{4\baselineskip}{\color{navy}\bfseries #1.}\enspace}
\newcommand{\example}[1]{\par\needspace{7\baselineskip}{\color{navy}\bfseries Разобранный пример.} #1\par}
\newcommand{\warning}[1]{\par\needspace{3\baselineskip}{\color{muted}\bfseries Типичная ошибка.} #1\par}
\newcommand{\selfcheck}[1]{\par\needspace{3\baselineskip}{\color{muted}\bfseries Проверка себя.} #1\par}
\newcommand{\minitrain}{\par\needspace{6\baselineskip}\vspace{1mm}{\color{navy}\bfseries Три мини-упражнения}\par\vspace{-1mm}}
\newcommand{\mini}[2]{\par\noindent\hspace*{2mm}\textbf{#1)}\enspace #2\par}
\newcommand{\task}[2]{\par\needspace{2\baselineskip}\noindent\makebox[8mm][l]{\color{navy}\bfseries #1.}#2\par\vspace{2mm}}
\newcommand{\training}[2]{\clearpage\section*{Тренировочный блок #1}\textcolor{muted}{#2}\par\vspace{3mm}\rule{\linewidth}{.4pt}\vspace{2mm}}
\newcommand{\answers}[2]{\clearpage\section*{Ответы к тренировочному блоку #1}\vspace{3mm}#2}
\newcommand{\anstable}[1]{\begin{center}\small\begin{tabularx}{\linewidth}{>{\centering\arraybackslash}p{9mm}X>{\centering\arraybackslash}p{9mm}X}\toprule\textbf{№}&\textbf{Ответ}&\textbf{№}&\textbf{Ответ}\\\midrule #1\\\bottomrule\end{tabularx}\end{center}}
\tikzset{
 st/.style={ellipse,draw=navy,line width=.8pt,fill=pale,align=center,minimum height=10mm,text width=43mm,font=\small},
 pr/.style={rectangle,rounded corners=2mm,draw=slate,line width=.8pt,align=center,minimum height=12mm,text width=68mm,font=\small},
 de/.style={diamond,aspect=2.45,draw=navy,line width=.8pt,align=center,text width=33mm,inner sep=1.7mm,font=\small},
 ar/.style={-{Stealth[length=2.4mm,width=1.7mm]},line width=.8pt,draw=navy},
 ann/.style={font=\footnotesize\bfseries,inner sep=1.5pt,color=navy}
}
\newcommand{\schemetitle}[1]{\clearpage\thispagestyle{fancy}\section*{Алгоритм: #1}\vspace{12mm}}

\begin{document}
\begin{titlepage}
\thispagestyle{empty}
\vspace*{23mm}
{\color{navy}\large УЧЕБНО-МЕТОДИЧЕСКОЕ ПОСОБИЕ}\par
\vspace{12mm}
{\fontsize{28}{35}\selectfont\bfseries Простые и\\[2mm] составные числа}\par
\vspace{8mm}
{\Large Задачи на целые числа. Подготовка к ЕГЭ\par}
\vspace{17mm}
{\color{muted}\large Определения, доказательства, методы,\\[2mm] тренировки и итоговая контрольная}\par
\vspace{15mm}
\begin{tikzpicture}[x=1cm,y=1cm]
\draw[navy,line width=.95pt] (0,0) .. controls (2.8,.9) and (6,-.65) .. (8.4,.14) .. controls (10.5,.77) and (12.8,.28) .. (15.6,.48);
\draw[slate,line width=.35pt] (0,-.36) -- (15.6,-.36);
\end{tikzpicture}
\vfill
\end{titlepage}

\clearpage
\section*{Карта темы}
\textbf{Главный вопрос:} как определить, является ли натуральное число простым, доказать составность выражения и использовать простые множители при решении задач с целыми числами?
\begin{center}
\begin{tikzpicture}[node distance=8mm]
\node[pr,text width=112mm] (a) {\textbf{1. Узнаём}\quad простое, составное, проверка делителей до $\sqrt n$};
\node[pr,text width=112mm,below=of a] (b) {\textbf{2. Доказываем составность}\quad разложение на множители и делимость};
\node[pr,text width=112mm,below=of b] (c) {\textbf{3. Исследуем структуру}\quad простые множители, степени, число делителей};
\node[pr,text width=112mm,below=of c] (d) {\textbf{4. Рассуждаем и конструируем}\quad простые числа, оценки, противоречие};
\draw[ar](a)--(b);\draw[ar](b)--(c);\draw[ar](c)--(d);
\end{tikzpicture}
\end{center}
\concept{Обозначения} $\mathbb N=\{1,2,3,\ldots\}$; запись $d\mid n$ означает, что $n$ делится на $d$; $d\nmid n$ --- что не делится. Число натуральных делителей $n$ обозначим $\tau(n)$, сумму этих делителей --- $\sigma(n)$.
\concept{План обучения} Сначала \emph{распознаём} числа, затем \emph{выбираем} проверку, \emph{конструируем} множители, \emph{сопоставляем} случаи и \emph{доказываем}, что иных вариантов нет.
\concept{Важное различие} Для утверждения «существует» достаточно примера. Для «невозможно», «при всех» и «найдите все» необходимо полное рассуждение.

% BLOCK 1
\blocktitle{1}{Понятие простого числа и проверка простоты}
\subsection{1.1. Простое, составное и особое число 1}
\concept{Определение} Натуральное число $p>1$ называют \emph{простым}, если у него ровно два различных натуральных делителя: $1$ и $p$. Натуральное число $n>1$ называют \emph{составным}, если у него есть делитель $d$ с $1<d<n$.
\concept{Сравнение} $13$ --- простое (делители $1,13$). $15$ --- составное ($15=3\cdot5$). Число $1$ имеет единственный делитель и потому \textbf{не является ни простым, ни составным}. Число $2$ --- единственное чётное простое число.
\concept{Главный критерий составности} Для $n>1$:
\[
 n\text{ составное}\quad\Longleftrightarrow\quad n=ab,\qquad a,b\in\mathbb N,\quad a>1,\ b>1.
\]
В одну сторону берём нетривиальный делитель $a$, а второй множитель равен $b=n/a$. В обратную сторону $a$ --- делитель, отличный от $1$ и $n$.
\example{Сравним числа $29$ и $35$. У $35$ сразу находится разложение $35=5\cdot7$, поэтому оно составное. Для $29$ проверим возможные делители $2,3,5$: ни один не подходит; почему этого достаточно, объясним в следующем подразделе. Следовательно, $29$ простое.}
\warning{Число с большим количеством цифр не обязательно составное; число с малым количеством цифр не обязательно простое. Проверять нужно делимость, а не внешний вид записи.}
\minitrain
\mini{1}{К какому виду относится $1$? (\textbf{Ни к простым, ни к составным}.)}
\mini{2}{Число $17$ простое или составное? (\textbf{Простое}.)}
\mini{3}{Докажите составность $49$. ($49=7\cdot7$.)}

\subsection{1.2. Почему достаточно делителей до $\sqrt n$}
\concept{Теорема} Если натуральное число $n>1$ составное, то у него есть делитель $d$, удовлетворяющий
\[
 1<d\le\sqrt n.
\]
\textbf{Доказательство.} Запишем $n=ab$, где $a,b>1$. Если бы и $a>\sqrt n$, и $b>\sqrt n$, то $ab>n$, что невозможно. Значит, хотя бы один множитель не превосходит $\sqrt n$.
\concept{Следствие} Для проверки простоты $n\ge2$ достаточно проверить делимость $n$ на \emph{простые числа} $p\le\sqrt n$. Если ни одно не делит $n$, число простое. Составной проверяемый делитель отдельно не нужен: у него самого есть простой делитель.
\example{Проверим $97$. Поскольку $9^2<97<10^2$, достаточно простых $2,3,5,7$. Число $97$ нечётно; сумма цифр $16$ не делится на $3$; последняя цифра не $0$ и не $5$; $97=7\cdot13+6$. Делителей не найдено, значит, $97$ простое. Для сравнения, $221=13\cdot17$: число составное, а $13<\sqrt{221}$.}
\selfcheck{Проверяются только делители, не превышающие $\sqrt n$, включая границу. Так, для $121$ обязательно проверить $11$: $121=11^2$.}
\minitrain
\mini{1}{До какого простого делителя проверяем $53$? (\textbf{До $7$ включительно}.)}
\mini{2}{Просто ли $101$? (\textbf{Да}: не делится на $2,3,5,7$.)}
\mini{3}{Составно ли $143$? (\textbf{Да}: $143=11\cdot13$.)}

\schemetitle{проверка простоты натурального числа}
\begin{center}
\begin{tikzpicture}[node distance=12mm]
\node[st] (s) {Дано натуральное $n$};
\node[de,below=15mm of s] (one) {$n=1$?};
\node[pr,text width=28mm,right=4mm of one] (neither) {Ни простое, ни составное};
\node[pr,below=18mm of one] (init) {Положите $d=2$};
\node[de,below=17mm of init] (lim) {$d^2>n$?};
\node[pr,text width=32mm,left=43mm of lim] (prime) {$n$ простое};
\node[de,below=20mm of lim] (div) {$d\mid n$?};
\node[pr,text width=36mm,left=43mm of div] (comp) {$n$ составное};
\node[pr,text width=36mm,below=18mm of div] (inc) {Замените $d$ на $d+1$};
\draw[ar](s)--(one);
\draw[ar](one.east)--node[ann,above]{да}(neither.west);
\draw[ar](one.south)--node[ann,right]{нет}(init.north);
\draw[ar](init)--(lim);
\draw[ar](lim.west)--node[ann,above]{да}(prime.east);
\draw[ar](lim.south)--node[ann,right]{нет}(div.north);
\draw[ar](div.west)--node[ann,above]{да}(comp.east);
\draw[ar](div.south)--node[ann,right]{нет}(inc.north);
\draw[ar](inc.east)--++(35mm,0)|-([yshift=0mm]lim.east);
\end{tikzpicture}
\end{center}

\training{1}{От распознавания числа --- к доказательству простоты и составности.}
\task{1}{Классифицируйте числа $1$, $2$, $9$, $17$, $21$.}
\task{2}{Докажите, что $91$ --- составное число.}
\task{3}{Найдите наименьшее простое число, большее $20$.}
\task{4}{Определите, просто ли число $97$.}
\task{5}{Проверьте, является ли число $119$ составным.}
\task{6}{Сколько простых чисел находится среди натуральных чисел от $10$ до $20$ включительно?}
\task{7}{Найдите все пары \emph{последовательных натуральных чисел}, оба числа в которых простые.}
\task{8}{Перечислите все простые числа в промежутке $40\le n\le50$.}
\task{9}{Найдите все натуральные $n$, при которых числа $n$, $n+2$, $n+4$ одновременно простые.}
\task{10}{Можно ли разместить $113$ одинаковых фишек в прямоугольник с целым числом строк и столбцов, если строк и столбцов должно быть не меньше двух? Обоснуйте.}

\answers{1}{\anstable{
1&$1$ --- ни то ни другое; $2,17$ --- простые; $9,21$ --- составные&6&$4$ (числа $11,13,17,19$)\\
2&$91=7\cdot13$&7&Только $(2,3)$\\
3&$23$&8&$41,43,47$\\
4&Да, простое&9&$n=3$\\
5&Да, $119=7\cdot17$&10&Нет: $113$ --- простое
}}

% BLOCK 2
\blocktitle{2}{Как доказывают составность выражений}
\subsection{2.1. Разложение на множители}
\concept{Цель} Не вычислять большое значение выражения, а представить его как произведение \emph{двух натуральных множителей, каждый больше $1$}.
\concept{Методы} Общий множитель, группировка, разность квадратов и разложение квадратного трёхчлена. При проверке последнего шага следует отдельно убедиться, что оба множителя больше единицы.
\example{Докажем, что при любом натуральном $n$ число $8n^2+10n+3$ составное. Подбираем множители и проверяем раскрытием скобок:
\[
8n^2+10n+3=(2n+1)(4n+3).
\]
При $n\ge1$ имеем $2n+1\ge3$ и $4n+3\ge7$. Оба множителя --- натуральные числа больше $1$. Значит, выражение составное при каждом $n\in\mathbb N$. Для проверки при $n=1$ получаем $21=3\cdot7$.}
\concept{Полезное сравнение} Число $n^2-1=(n-1)(n+1)$ составное только при $n\ge3$: если $n=2$, получаем $3=1\cdot3$, а единица не делает число составным.
\warning{Запись $N=1\cdot N$ не является доказательством составности. И запись $N=ab$ недостаточна, если один из множителей может быть равен $1$ или неположителен.}
\minitrain
\mini{1}{Разложите $9k^2-1$. ($\boldsymbol{(3k-1)(3k+1)}$.)}
\mini{2}{Составно ли $n^2+5n+6$ при $n\ge1$? (\textbf{Да}: $(n+2)(n+3)$.)}
\mini{3}{При каких натуральных $n$ число $n^2-1$ составное? ($\boldsymbol{n\ge3}$.)}

\subsection{2.2. Чётность, кратность и большие степени}
\concept{Новый подход} Иногда нет необходимости полностью раскладывать выражение. Достаточно показать, что оно делится на некоторое $m>1$ и строго больше $m$: тогда оно составное.
\example{Докажем, что $n(n+1)+2$ составно для всех $n\ge1$. Два соседних числа $n,n+1$ имеют разную чётность, поэтому $n(n+1)$ чётно. Следовательно, $n(n+1)+2$ чётно и не меньше $4$. Значит, его делит $2$, причём оно больше $2$, и потому составно.}
\concept{Факторизация степеней} Для любого натурального $n\ge2$:
\[
2^{2n}-1=(2^n-1)(2^n+1).
\]
Оба множителя больше $1$, так что $2^{2n}-1$ составно. При $n=1$ это число равно $3$, то есть ограничение $n\ge2$ существенно.
\example{Не вычисляя $2^{20}$, докажем составность $2^{20}-1$:
\[
2^{20}-1=(2^{10}-1)(2^{10}+1)=1023\cdot1025.
\]
Оба множителя больше $1$, проверка выполняется обратным раскрытием скобок.}
\minitrain
\mini{1}{Составно ли $2n+4$ при $n\ge1$? (\textbf{Да}: чётно и не меньше $6$.)}
\mini{2}{Разложите $3^8-1$ как разность квадратов. ($\boldsymbol{(3^4-1)(3^4+1)}$.)}
\mini{3}{Делится ли $n^3-n$ на $6$ при любом целом $n$? (\textbf{Да}.)}

\schemetitle{доказательство составности}
\begin{center}
\begin{tikzpicture}[node distance=13mm]
\node[st] (s) {Дано натуральное выражение $N>1$};
\node[pr,below=of s] (f) {Попробуйте разложить $N$ на множители};
\node[de,below=17mm of f] (q) {Есть $N=ab$, $a,b>1$?};
\node[pr,text width=35mm,left=45mm of q] (yes) {$N$ составное};
\node[pr,below=18mm of q] (g) {Найдите делитель $m>1$ по чётности или кратности};
\node[de,below=17mm of g] (t) {$m\mid N$ и $N>m$?};
\node[pr,text width=35mm,left=45mm of t] (yes2) {$N$ составное};
\node[pr,text width=39mm,below=18mm of t] (no) {Метод не дал вывода; изучите другие делители};
\draw[ar](s)--(f);\draw[ar](f)--(q);
\draw[ar](q.west)--node[ann,above]{да}(yes.east);
\draw[ar](q.south)--node[ann,right]{нет}(g.north);
\draw[ar](g)--(t);
\draw[ar](t.west)--node[ann,above]{да}(yes2.east);
\draw[ar](t.south)--node[ann,right]{нет}(no.north);
\end{tikzpicture}
\end{center}

\training{2}{Доказательства для всех натуральных $n$ и анализ границ применимости.}
\task{1}{Разложите на множители $49k^2-1$.}
\task{2}{Докажите составность $n(n+1)+4$ для любого натурального $n$.}
\task{3}{Докажите, что $(3n+1)^2-1$ составно при всех $n\ge1$.}
\task{4}{Докажите составность $4^{10}-1$, не вычисляя само число.}
\task{5}{Докажите составность $n^2+7n+12$ при $n\ge1$.}
\task{6}{Докажите, что $6n^2+11n+3$ составно для любого $n\ge1$.}
\task{7}{Докажите, что $2^{4n}-1$ кратно $15$ для любого $n\ge1$.}
\task{8}{Докажите делимость $n(n+1)(n+2)$ на $6$ для любого целого $n$.}
\task{9}{Докажите составность $n^4+4$ для каждого натурального $n\ge2$. Укажите, что происходит при $n=1$.}
\task{10}{На электронном табло появляется натуральное число $N=t^2+t+2$, где $t\ge1$. Может ли оно быть простым? Докажите ответ.}

\answers{2}{\anstable{
1&$(7k-1)(7k+1)$&6&$(3n+1)(2n+3)$\\
2&Чётное число $\ge6$&7&Да: $16^n-1$ кратно $15$\\
3&$3n(3n+2)$&8&Да: среди трёх чисел есть кратное $3$ и чётное\\
4&$(2^{10}-1)(2^{10}+1)$&9&$(n^2-2n+2)(n^2+2n+2)$; при $n=1$ простое $5$\\
5&$(n+3)(n+4)$&10&Нет: чётное число $\ge4$
}}

% BLOCK 3
\blocktitle{3}{Простые множители и устройство натурального числа}
\subsection{3.1. Каноническое разложение на простые множители}
\concept{Основная теорема арифметики} Каждое натуральное число $n>1$ можно представить в виде произведения простых множителей; с точностью до их порядка такое представление единственно. Группируя одинаковые множители, получаем
\[
n=p_1^{\alpha_1}p_2^{\alpha_2}\cdots p_r^{\alpha_r},\qquad
\alpha_i\in\mathbb N,\quad p_i\ne p_j\ (i\ne j).
\]
Это \emph{каноническое разложение} натурального числа.
\concept{Почему существует разложение} Если $n$ простое, разложение найдено. Если оно составное, $n=ab$ с $1<a,b<n$; применяем то же рассуждение к $a$ и $b$. При каждом разбиении множители уменьшаются, поэтому процесс заканчивается простыми числами.
\concept{Почему оно единственно} Используется \emph{лемма Евклида}: если простое $p$ делит $ab$, то $p$ делит хотя бы один из множителей. Поэтому простой множитель из одного разложения обязательно присутствует в другом; его можно сократить и повторить рассуждение.
\example{Разложим $756$:
\[
756=4\cdot189=2^2\cdot(27\cdot7)=2^2\cdot3^3\cdot7.
\]
Проверим: $4\cdot27\cdot7=756$. У числа три различных простых делителя: $2,3,7$.}
\warning{Числа $2^3\cdot3$ и $2\cdot3^3$ не совпадают: показатели степеней нельзя переставлять между разными простыми числами.}
\minitrain
\mini{1}{Разложите $84$ на простые множители. ($\boldsymbol{2^2\cdot3\cdot7}$.)}
\mini{2}{Сколько различных простых делителей у $90$? ($\boldsymbol{3}$: $2,3,5$.)}
\mini{3}{Является ли $121$ степенью простого числа? (\textbf{Да}: $11^2$.)}

\subsection{3.2. Степени простых чисел, количество и сумма делителей}
\concept{Примарные числа} Натуральное число вида $p^k$, где $p$ --- простое и $k\ge1$, называют \emph{степенью простого числа} (примарным числом). Например, $27=3^3$ примарно, а $12=2^2\cdot3$ --- нет. Простое число --- частный случай при $k=1$.
\concept{Количество делителей} Если $n=p_1^{\alpha_1}\cdots p_r^{\alpha_r}$, то каждый натуральный делитель имеет вид $p_1^{\beta_1}\cdots p_r^{\beta_r}$, где $0\le\beta_i\le\alpha_i$. Поэтому
\[
\boxed{\tau(n)=(\alpha_1+1)(\alpha_2+1)\cdots(\alpha_r+1).}
\]
\concept{Сумма делителей} Раскрывая произведение конечных сумм, получаем
\[
\boxed{\sigma(n)=(1+p_1+\cdots+p_1^{\alpha_1})\cdots
(1+p_r+\cdots+p_r^{\alpha_r}).}
\]
\example{Для $72=2^3\cdot3^2$:
\[
\tau(72)=(3+1)(2+1)=12,\qquad
\sigma(72)=(1+2+4+8)(1+3+9)=15\cdot13=195.
\]
Следовательно, у $72$ ровно $12$ натуральных делителей, а сумма всех делителей равна $195$.}
\concept{Обратная задача} Ровно три натуральных делителя имеют только числа $p^2$ с простым $p$. Ровно четыре делителя имеют числа $p^3$ либо $pq$ для различных простых $p,q$: это следует из разложений $4=4$ и $4=2\cdot2$ для формулы $\tau$.
\minitrain
\mini{1}{Найдите $\tau(48)$. ($\boldsymbol{10}$: $48=2^4\cdot3$.)}
\mini{2}{Найдите наименьшее натуральное число с тремя делителями. ($\boldsymbol{4}$.)}
\mini{3}{Найдите $\sigma(12)$. ($\boldsymbol{28}$.)}

\schemetitle{исследование числа по простым множителям}
\begin{center}
\begin{tikzpicture}[node distance=13mm]
\node[st] (s) {Натуральное $n>1$};
\node[pr,below=of s] (a) {Разложите $n$ в произведение простых степеней};
\node[pr,below=of a] (b) {Запишите показатели $\alpha_1,\ldots,\alpha_r$};
\node[de,below=18mm of b] (q) {Нужно найти $\tau(n)$?};
\node[pr,text width=43mm,below=28mm of q,xshift=-49mm] (l) {Количество: $\tau(n)=\prod(\alpha_i+1)$};
\node[pr,text width=47mm,below=28mm of q,xshift=49mm] (r) {Сумму: $\sigma(n)=\prod(1+p_i+\cdots+p_i^{\alpha_i})$};
\node[st,text width=53mm,below=18mm of q,yshift=-60mm] (end) {Проверьте результат на небольшом примере};
\draw[ar](s)--(a);\draw[ar](a)--(b);\draw[ar](b)--(q);
\draw[ar](q.west)-|node[ann,near start,above]{да}(l.north);\draw[ar](q.east)-|node[ann,near start,above]{нет}(r.north);
\draw[ar](l.south)|-(end.west);\draw[ar](r.south)|-(end.east);
\end{tikzpicture}
\end{center}

\training{3}{От канонического разложения --- к классификации чисел по количеству делителей.}
\task{1}{Разложите число $180$ на простые множители.}
\task{2}{Разложите число $98$ на простые множители.}
\task{3}{Проверьте, является ли $169$ простым числом.}
\task{4}{Сколько различных простых делителей имеет $840$?}
\task{5}{Найдите количество натуральных делителей числа $72$.}
\task{6}{Найдите наименьшее натуральное число с ровно тремя натуральными делителями.}
\task{7}{Найдите все числа $n<50$, у которых ровно три натуральных делителя.}
\task{8}{Найдите все натуральные числа $n<25$, имеющие ровно четыре натуральных делителя.}
\task{9}{Найдите наименьшее натуральное число, имеющее ровно шесть делителей и только один различный простой делитель.}
\task{10}{Найдите наименьшее натуральное число $N$, которое делится на $2$, $3$ и $5$ и имеет ровно $12$ натуральных делителей.}

\answers{3}{\anstable{
1&$2^2\cdot3^2\cdot5$&6&$4$\\
2&$2\cdot7^2$&7&$4,9,25,49$\\
3&Нет: $169=13^2$&8&$6,8,10,14,15,21,22$\\
4&$4$ ($2,3,5,7$)&9&$2^5=32$\\
5&$12$&10&$60=2^2\cdot3\cdot5$
}}

% BLOCK 4
\blocktitle{4}{Доказательные задачи о простых числах}
\subsection{4.1. Бесконечность простых чисел и длинные ряды составных}
\concept{Теорема Евклида} Простых натуральных чисел бесконечно много.
\textbf{Доказательство.} Предположим, что все простые числа перечислены: $p_1,\ldots,p_k$. Рассмотрим
\[
N=p_1p_2\cdots p_k+1>1.
\]
Число $N$ имеет хотя бы один простой делитель $q$. Однако при делении $N$ на каждый $p_i$ остаток равен $1$, поэтому $q$ не совпадает ни с одним $p_i$. Получено противоречие с полнотой списка.
\warning{Число «произведение всех выбранных простых плюс $1$» \emph{не обязано} быть простым. Например,
\[
2\cdot3\cdot5\cdot7\cdot11\cdot13+1=30031=59\cdot509.
\]
Доказательство говорит о \emph{новом простом делителе}, а не о простоте всего числа.}
\concept{Сколько угодно подряд составных} Для любого $m\ge2$ числа
\[
(m+1)!+2,\ (m+1)!+3,\ \ldots,\ (m+1)!+(m+1)
\]
составные: каждое число с добавкой $k$ делится на $k$ и больше $k$. Таких последовательных чисел ровно $m$.
\example{Чтобы получить шесть составных чисел подряд, берём $7!+2,7!+3,\ldots,7!+7$, то есть $5042,5043,\ldots,5047$. Каждое делится соответственно на $2,3,\ldots,7$ и строго больше делителя.}
\minitrain
\mini{1}{Найдите $2\cdot3\cdot5+1$. Простое ли число? ($\boldsymbol{31}$, \textbf{да}.)}
\mini{2}{Сколько составных подряд даёт набор $6!+2,\ldots,6!+6$? ($\boldsymbol{5}$.)}
\mini{3}{Верно ли, что $p_1\cdots p_k+1$ всегда простое? (\textbf{Нет}, например $30031$.)}

\subsection{4.2. Простые числа и ограничения на остатки}
\concept{Смысл} Условие простоты часто запрещает число делиться на малое простое. Например, простое $p>3$ нечётно и не делится на $3$. Поэтому при делении на $6$ оно может иметь только остаток $1$ или $5$, то есть
\[
p=6k\pm1\qquad(k\in\mathbb Z).
\]
Обратное утверждение неверно: $25=6\cdot4+1$, но $25$ составное.
\example{Докажем, что для любого простого $p>3$ число $p^2-1$ делится на $24$. Запишем
\[
p^2-1=(p-1)(p+1).
\]
Оба числа $p-1$ и $p+1$ чётные; одно из них кратно $4$, значит, их произведение делится на $8$. Поскольку $p$ не кратно $3$, одно из соседей $p-1,p+1$ кратно $3$. Следовательно, произведение делится одновременно на взаимно простые $8$ и $3$, то есть на $24$. Проверка для $p=5$: $25-1=24$.}
\concept{Тройка простых} Если $p,p+2,p+4$ --- простые, то ровно одно из них делится на $3$. Оно должно равняться $3$, а значит $p=3$ и получается тройка $3,5,7$. Иных троек нет.
\warning{Из записи $p=6k\pm1$ не следует простота $p$. Это \emph{необходимый}, но не достаточный признак.}
\minitrain
\mini{1}{Какой остаток при делении на $6$ может давать простое $p>3$? ($\boldsymbol{1\text{ или }5}$.)}
\mini{2}{Делится ли $p^2-1$ на $8$ для простого $p>2$? (\textbf{Да}.)}
\mini{3}{Найдите все простые $p$, для которых $p,p+2,p+4$ простые. ($\boldsymbol{p=3}$.)}

\schemetitle{ограничение возможных простых чисел}
\begin{center}
\begin{tikzpicture}[node distance=13mm]
\node[st] (s) {Дано условие на простое $p$};
\node[pr,below=of s] (a) {Исключите малые простые $2,3,5$ или проверьте их отдельно};
\node[pr,below=of a] (b) {Рассмотрите остатки по подходящему модулю};
\node[de,below=18mm of b] (q) {Есть кратное малому простому?};
\node[pr,text width=42mm,below=29mm of q,xshift=-49mm] (l) {Противоречие, если число больше своего делителя};
\node[pr,text width=42mm,below=29mm of q,xshift=49mm] (r) {Оставьте возможные варианты};
\node[st,text width=51mm,below=18mm of q,yshift=-61mm] (end) {Сформулируйте вывод и проверьте кандидатов};
\draw[ar](s)--(a);\draw[ar](a)--(b);\draw[ar](b)--(q);
\draw[ar](q.west)-|(l.north);\draw[ar](q.east)-|(r.north);
\node[ann] at ([xshift=-11mm,yshift=2mm]q.west) {да};
\node[ann] at ([xshift=11mm,yshift=2mm]q.east) {нет};
\draw[ar](l.south)|-(end.west);\draw[ar](r.south)|-(end.east);
\end{tikzpicture}
\end{center}

\training{4}{Противоречие, простые делители, последовательности и конструктивные примеры.}
\task{1}{Вычислите $2\cdot3\cdot5\cdot7+1$ и определите, простое ли полученное число.}
\task{2}{Объясните, почему числа $6!+2,6!+3,\ldots,6!+6$ все составные.}
\task{3}{Докажите, что любые два последовательных натуральных числа взаимно просты.}
\task{4}{Найдите все простые $p$, для которых $p,p+4,p+8$ --- простые числа.}
\task{5}{Докажите, что при любом простом $p>3$ число $p^3-p$ кратно $24$.}
\task{6}{Докажите, что среди любых пяти последовательных целых чисел есть число, кратное $5$.}
\task{7}{Докажите, что всякое простое $p>3$ имеет вид $6k-1$ или $6k+1$, где $k$ --- целое число.}
\task{8}{Укажите семь последовательных составных натуральных чисел, не проверяя их делимость по отдельности.}
\task{9}{Докажите, что для каждого простого $p>3$ число $p^4+23$ делится на $24$.}
\task{10}{Постройте десять последовательных составных натуральных чисел так, чтобы число, стоящее \emph{перед} ними, делилось на $1000$. Укажите первое число последовательности и обоснуйте построение.}

\answers{4}{\anstable{
1&$211$, простое&6&Да (пять разных остатков по модулю $5$)\\
2&Каждое $6!+k$ кратно $k$ ($2\le k\le6$)&7&$p\equiv1$ или $5\pmod6$\\
3&Общий делитель делит разность $1$&8&$8!+2,\ldots,8!+8$\\
4&$p=3$&9&Да: $p^2-1$ кратно $24$\\
5&$p(p^2-1)$ кратно $24$&10&$B+1,\ldots,B+10$, где $B=11!\cdot1000$
}}

\selfcheck{К заданию 10: $B=39\,916\,800\,000$ кратно $1000$; $B+1=13\cdot3\,070\,523\,077$ составное, а каждое $B+k$ при $2\le k\le10$ делится на $k$ и больше $k$.}

% FINAL
\clearpage
\section*{Итоговая обзорная контрольная работа}
\textcolor{muted}{10 заданий: 1--3 --- базовые, 4--7 --- средние, 8--9 --- повышенные, 10 --- комплексное.}
\par\vspace{3mm}\rule{\linewidth}{.4pt}\vspace{2mm}
\task{1}{Классифицируйте числа $1,29,39,51$ как простые, составные или не относящиеся ни к одной из этих групп.}
\task{2}{Разложите число $126$ на простые множители.}
\task{3}{Проверьте, является ли число $137$ простым. Обоснуйте полноту проверки.}
\task{4}{Докажите, что $(n+1)(n+2)$ составно при любом натуральном $n$.}
\task{5}{Найдите число натуральных делителей $180$.}
\task{6}{Найдите все натуральные числа, не превосходящие $100$, имеющие ровно три натуральных делителя.}
\task{7}{Докажите, что число $n^2+n+2$ составное при любом натуральном $n$.}
\task{8}{Найдите все простые числа $p$, для которых $p,p+2,p+4$ одновременно простые.}
\task{9}{Докажите, что для каждого простого $p\ge5$ число $p^4-1$ делится на $48$.}
\task{10}{Найдите наименьшее натуральное число $N$, которое кратно $30$ и имеет ровно $24$ натуральных делителя. Докажите минимальность.}

\clearpage
\section*{Ответы к итоговой обзорной контрольной работе}
\anstable{
1&$1$ --- ни то ни другое; $29$ --- простое; $39,51$ --- составные&6&$4,9,25,49$\\
2&$126=2\cdot3^2\cdot7$&7&Составное, так как чётное и $\ge4$\\
3&$137$ --- простое&8&$p=3$\\
4&Да, оба множителя больше $1$&9&Да: $(p^2-1)(p^2+1)$ кратно $24\cdot2$\\
5&$18$&10&$360=2^3\cdot3^2\cdot5$
}
\vspace{8mm}
\section*{Памятка: что проверять в задачах на простоту}
\begin{itemize}
\item \textbf{Определения:} простое число имеет ровно два натуральных делителя; $1$ --- ни простое, ни составное.
\item \textbf{Проверка простоты:} искать простые делители не выше $\sqrt n$.
\item \textbf{Доказательство составности:} предъявить два натуральных множителя больше $1$ либо нетривиальный делитель.
\item \textbf{Структура числа:} каноническое разложение задаёт число делителей по формуле $\tau(n)=\prod(\alpha_i+1)$.
\item \textbf{Сложные задачи:} использовать малые простые делители, остатки, ограничения, противоречие и конструкции вида $m!+k$.
\item \textbf{Полнота:} найденные кандидаты проверить; отсутствие остальных обосновать.
\end{itemize}
\end{document}